\section{JUSTIFICACIÓN}

\subsection{Justificación Científica y Tecnológica}

La adopción de esquemas criptográficos de clave pública basados en la dificultad de factorización de enteros, como RSA, en sistemas de votación electrónica desarrollados en el contexto universitario boliviano (Churata Sonco, 2020; Fernández Tristán, 2021), no ha sido examinada frente a la disponibilidad de estándares criptográficos post-cuánticos. El National Institute of Standards and Technology (2024) publicó los estándares FIPS 203 y FIPS 204, cuya seguridad se sustenta en problemas matemáticos distintos de los vulnerables a algoritmos cuánticos como el de Shor (Fortinet, s.f.). La investigación se relaciona con el estado del arte en verificación formal de protocolos de votación electrónica, en el cual Baloglu et al. (2021) evidenciaron que las propiedades de seguridad de un protocolo dependen tanto de sus primitivas criptográficas como de su diseño formal. El estudio se sitúa en el vacío identificado entre la disponibilidad de dichos estándares y su aplicación específica al diseño y verificación formal de protocolos de votación electrónica en el contexto universitario boliviano.

\subsection{Justificación Metodológica}

El problema de investigación plantea una comparación entre dos tipos de protocolo criptográfico —post-cuántico y clásico— respecto de variables de integridad, confidencialidad y rendimiento, lo que corresponde a un diseño que combina análisis formal y medición empírica. La verificación formal mediante Tamarin Prover permite evaluar propiedades de integridad, confidencialidad y autenticidad de forma sistemática, siguiendo un enfoque metodológico ya aplicado en el análisis de protocolos de votación electrónica como Helios y Belenios (Baloglu et al., 2021). La medición de indicadores de rendimiento mediante implementación y simulación permite obtener valores comparables entre el protocolo post-cuántico propuesto y un protocolo equivalente basado en RSA/ECDSA. El procedimiento descrito, al basarse en herramientas y estándares documentados públicamente, resulta reproducible en investigaciones posteriores que aborden otros protocolos criptográficos o contextos de aplicación distintos al proceso de votación electrónica.

\subsection{Justificación Social y Académica}

El conocimiento generado por la investigación resulta pertinente para la comunidad académica de la Carrera de Informática de la Universidad Mayor de San Andrés, en la medida en que examina un problema no abordado en los antecedentes locales identificados (Churata Sonco, 2020; Fernández Tristán, 2021), los cuales se limitan a esquemas criptográficos clásicos. El estudio también se relaciona con investigadores y desarrolladores que trabajan en sistemas de votación electrónica o en la aplicación de criptografía post-cuántica a protocolos con requerimientos de integridad y confidencialidad. Los resultados de la verificación formal y de la comparación de rendimiento entre esquemas post-cuánticos y clásicos podrían constituir una referencia para investigaciones futuras que analicen la transición de otros protocolos de comunicación hacia primitivas criptográficas resistentes a la computación cuántica, dentro del área de la seguridad informática.
